\clearpage
\addcontentsline{toc}{section}{Список литературы}
\begin{thebibliography}{9}
\bibitem{course} Кухтичев А.\,А. Материалы курса «Информационный поиск»,
МАИ, осень 2025: задания, лекции и шаблон отчёта.
Репозиторий MAI-IR: условия лабораторных работ,
лекционные материалы и домашние задания лекций 1, 2, 4.
\bibitem{manning} Маннинг К.\,Д., Рагхаван П., Шютце Х.
Введение в информационный поиск. М.: Вильямс, 2011. 528 с.
\bibitem{dataset} Common Pile. Stack Exchange: архивный корпус.
\url{https://huggingface.co/datasets/common-pile/stackexchange}.
Использованная ревизия указана в лабораторной №1.
\bibitem{wordnet} Princeton University. WordNet 3.0.
Словарные данные и исключения, включённые в исходники проекта;
лицензия сохранена вместе со словарём.
\bibitem{porter} Porter M.\,F. An algorithm for suffix stripping.
Program. 1980. Vol. 14, no. 3. P. 130--137.
\bibitem{implementation} Исходные тексты лабораторных работ №1--8,
10--13, 17. Собственные реализации индексатора, поисковых алгоритмов,
оценки качества и автоматических тестов; результаты экспериментов
по корпусу документации и обсуждений программирования.
\end{thebibliography}
